% Transcription — fonds Alexandre Grothendieck, shelfmark 109, batch 2 (pages 21-37).
% Established from the facsimile archives/batches/109/batch-02.pdf (a mirror of
% https://grothendieck.umontpellier.fr/109.pdf), per the skill
% transcribe-grothendieck. The batch file carries no cover sheet: PDF page k
% is page 20+k of the folder (the Montpellier footer prints « 21 » ... « 37 »),
% and the archivists' pencil numbers (bottom left, top right on page 28)
% agree on every leaf. The folder is thirty-seven pages; this is the second
% and last batch.
%
% Pass: Opus 5.5 (claude-opus-5-5), 2026-09-24 — first pass, unchecked against the pages by a human.
%
% Skipped pages: none. All seventeen pages carry \page markers.
%
% Summarised as another's text: all seventeen pages, 21-37. They are
% photocopies of one handwritten list, in English, in a single rounded hand
% that is not his: a native English speaker's page-by-page corrections to an
% English typescript of his, addressed to him in the second person (« do you
% mean », « I suggest », « can you check? »), sent in instalments — batch 1's
% page 20 carries « Sent 19/12/83 » and page 28 here « Sent 6/2/84 ». The
% writer is not named on these leaves and is not named here. The list is not
% a letter on the mathematics (nine entries in ten are English usage:
% contractions, « its »/« it's », « commutes with », spelling), so it is
% treated under the rule for other people's text — one French \note per
% page, no recomposition — rather than as a letter transcribed in full. The
% few entries that bear on the mathematics or the notation (a formula
% queried, a symbol missing on the copy, a notation asked for) are named in
% the notes, with the formula where the list writes one. This choice should
% be checked against what batch 1 does with pages 18 and 20, which open the
% same list. No mark in his hand was found on any of the seventeen pages;
% the small circles near the top and bottom right corners are marks of the
% copying and are not reproduced.
%
% What the batch is about. The list runs on from batch 1 (whose page 20
% stops at typescript p. 115) through typescript pp. 116-557, in order,
% one or two typescript pages per line: 21 (pp. 116-154), 22 (156-190),
% 23 (191-224), 24 (225-274), 25 (275?-302), 26 (302-334), 27 (335-354),
% 28 (335-371, with « Sent 6/2/84 » after p. 354), 29 (371-401), 30
% (402-425), 31 (426-468), 32 (468-482), 33 (483-494), 34 (495-514), 35
% (515-532), 36 (532-549), 37 (550-557, the list stopping a third of the
% way down). Page 27 is a copy of the first twenty lines of page 28 in an
% earlier state, before the date and the continuation were added. The
% typescript is not named on the leaves; the inventory group (« Autour de
% Pursuing stacks ») and what the entries quote — test categories, W,
% asphericity, j_A : Cat -> Â, homotopy intervals, the Kan condition,
% « formulary », Whitehead, Illusie, Breen, Lazard — point to Pursuing
% Stacks, and the notes say « sans doute » once and no more. Mathematical
% queries worth a look: p. 116 (K_h c Z_h), 224 (j_A), 234-235 (asphericity
% « a tautology as it stands »), 240 (W = (i_{M,i})^{-1}(W)), 246 (i_0 :
% N_0 -> N ?), 269 (pi_0, pi_1 undefined), 327 (alpha : Delta -> Delta/F),
% 428 (Â_M, not Â, abelian), 471 (group object pi_1(X_*) in U(k)), 494
% (Sym^j_k, Gamma^j_k, SGA 6), 509 (LH and E), 517 (V and V'), 538
% (H^i(T, Pi_i(X_*)) = 0).
%
% Readings to check first: the subscripts in the formula of p. 116 (page
% 21); « C ≅ A ? » at the head of page 25 (the relation sign is
% overwritten); the expression at typescript p. 485 (page 33); the
% interlinear words under the j_A formula of p. 224 (page 23).

\documentclass[11pt,a4paper]{article}
% Le préambule commun à toutes les transcriptions du fonds Grothendieck.
%
% Il tient en une idée : **l'incertitude se compose, elle ne se résout pas.**
% Une transcription qui lisse un mot douteux a détruit la seule chose pour
% laquelle on la faisait — un lecteur doit pouvoir savoir, à chaque endroit,
% ce qui était sur le feuillet et ce qui a été ajouté. Les six macros qui
% suivent sont donc l'appareil critique, et non de la décoration.
%
% Les mêmes commandes sont reconnues par `scripts/render.mjs`, qui produit la
% vue de lecture du site. Une macro ajoutée ici doit l'être aussi là-bas,
% faute de quoi le rendu échouera bruyamment — ce qui est voulu : un rendu
% silencieusement faux est pire qu'un rendu absent.

\ProvidesPackage{grothendieck}[2026/08/08 Transcriptions du fonds Grothendieck]

\usepackage{amsmath,amssymb,amsthm}
\usepackage{mathrsfs}
\usepackage{stmaryrd}
% Les diagrammes commutatifs sont partout dans ce fonds : c'est la langue
% même de la géométrie algébrique de Grothendieck, et une transcription qui
% les rendrait en prose ne transcrirait rien.
\usepackage{tikz-cd}
% Français par défaut — c'est la langue des feuillets — mais l'anglais est
% chargé aussi : la traduction et le résumé sont des documents anglais, et la
% typographie française (espaces avant les deux-points) y serait une faute.
% Ces documents basculent par \selectlanguage{english} en tête de corps.
\usepackage[english,french]{babel}
\usepackage{fontspec}
\usepackage{geometry}
\geometry{a4paper,margin=25mm,marginparwidth=28mm}
\usepackage{marginnote}
\usepackage{xcolor}
% ulem, not soul. `soul`'s \st and \ul are built on a character-by-character
% scan that XeLaTeX's font machinery breaks: the document fails with an
% undefined control sequence rather than degrading. ulem does the same two jobs
% and is Unicode-engine safe. `normalem` stops it hijacking \emph, which the
% transcriptions use for Grothendieck's own emphasis.
\usepackage[normalem]{ulem}
\usepackage{hyperref}
\hypersetup{colorlinks=true,linkcolor=black,urlcolor=black,pdfborder={0 0 0}}

% --- Métadonnées du lot -----------------------------------------------------
% Reprises telles quelles par le script de rendu, qui les affiche en tête de
% la vue de lecture. Elles fixent ce que le fichier prétend couvrir : sans
% elles, une transcription flotte sans point d'attache dans un fonds de seize
% mille feuillets.

\newcommand{\folder}[1]{\gdef\grothFolder{#1}}
\newcommand{\batch}[1]{\gdef\grothBatch{#1}}
\newcommand{\leaves}[2]{\gdef\grothFirst{#1}\gdef\grothLast{#2}}
\newcommand{\foldertitle}[1]{\gdef\grothFoldertitle{#1}}
\gdef\grothFolder{?}\gdef\grothBatch{?}\gdef\grothFirst{?}\gdef\grothLast{?}\gdef\grothFoldertitle{}

% --- Le résumé --------------------------------------------------------------
%
% Il ouvre la lecture modernisée, sous le titre « Résumé », et il est composé
% en petit. Ce n'est pas un ornement : c'est le seul passage du document qui
% ne soit pas des mathématiques, et un lecteur doit pouvoir le reconnaître
% avant de l'avoir lu — puis le sauter s'il connaît déjà le sujet, ou s'y
% arrêter s'il ne le connaît pas. Le filet qui le suit marque la reprise du
% corps.

\newenvironment{resume}
  {\small\setlength{\parskip}{0.45em}}
  {\par\medskip\hrule\medskip}

% --- Mots-clés ---------------------------------------------------------------
%
% En anglais, en clôture du résumé de la lecture modernisée : le vocabulaire
% moderne sous lequel chercher ce que les feuillets construisent. Le manifeste
% du site (`npm run manifest`) les extrait pour étiqueter la cote entière —
% cette ligne est la source unique des tags, il n'y a pas de fichier à côté.
\newcommand{\keywords}[1]{\par\smallskip\noindent
  {\footnotesize\textsc{Keywords} --- \textit{#1}}\par}

% --- L'appareil critique ----------------------------------------------------

% Le numéro de feuillet, en marge. C'est l'ancre qui permet de régler un
% désaccord sur une formule en regardant *un* feuillet plutôt qu'une chemise
% de six cents. Le site s'en sert aussi pour tourner les pages du fac-similé
% quand on fait défiler la transcription.
\newcommand{\leaf}[1]{%
  \marginnote{\footnotesize\color{gray!70}\textsf{#1}}%
  \label{leaf:#1}\ignorespaces}

% Illisible : on ne devine pas. Un mot inventé qui se lit comme les autres est
% le pire résultat possible.
\newcommand{\ill}{\textcolor{red!60!black}{[\,\ldots\,]}}

% Lecture douteuse : le mot est donné, et signalé comme incertain.
\newcommand{\uncertain}[1]{\uline{#1}}

% Ajout de l'éditeur — une lettre manquante, un mot sous-entendu.
\newcommand{\add}[1]{\textcolor{blue!50!black}{[#1]}}

% Biffé par Grothendieck lui-même. Ce qu'il a rayé fait partie du document :
% une rature dit souvent où la pensée a changé de direction.
\newcommand{\struck}[1]{\sout{#1}}

% Note du transcripteur, sur l'état matériel du feuillet.
\newcommand{\note}[1]{\par\smallskip{\small\color{gray!60!black}\textit{[#1]}}\par\smallskip}

% Note marginale de Grothendieck, distincte du corps du texte.
\newcommand{\marginal}[1]{\par\smallskip\noindent\hspace{1em}%
  {\small\color{gray!60!black}#1}\par\smallskip}

% --- Titre ------------------------------------------------------------------

\AtBeginDocument{%
  \begin{center}
    {\large\bfseries Fonds Alexandre Grothendieck --- cote n\textsuperscript{o}~\grothFolder}\\[2pt]
    {\itshape\grothFoldertitle}\\[4pt]
    {\small feuillets \grothFirst--\grothLast\ (lot \grothBatch)}\\[2pt]
    {\footnotesize\color{gray!60!black}Transcription établie d'après le fac-similé numérisé
     par l'Université de Montpellier.\\
     Les lectures incertaines sont soulignées, les illisibles notées [\,\ldots\,],
     les ajouts entre crochets.}
  \end{center}
  \vspace{1em}}

\endinput


\folder{109}
\batch{2}
\pages{21}{37}
\dating{1983-1984}
\watermark{Édition de démonstration}
\foldertitle{Tables [Corrections et préparation à l'édition] : notes manuscrites (s.d.), copies de notes manuscrites annotées (s.d.), copies de tapuscrit annoté (s.d.), copies de notes manuscrites (1983-1984)}

\begin{document}

\page{21}
\note{les pages 21 à 37 sont la copie d'une liste manuscrite en anglais,
d'une autre main que la sienne, qui lui est adressée : des corrections,
page par page et ligne par ligne, à un tapuscrit anglais de lui (sans doute
celui de \emph{Pursuing Stacks}, d'après le groupe de l'inventaire et ce
que citent les renvois). La liste commence dans le lot 1 du dossier ; elle
est ici résumée page à page, et non recomposée. On n'y a relevé aucune
marque de sa main.}

\note{p.\ 116 à 154 du tapuscrit. À la p.\ 116, une formule à corriger :
\uncertain{$K_h \subset K_n$} devrait être $K_h \subset Z_h$. Pour les
p.\ 114--115, l'auteur de la liste demande si l'annotation manuscrite,
qui renvoie au §\,51, ne vise pas plutôt le §\,52. À la p.\ 151,
une reformulation sur la catégorie but (Cat) de la formule (2), et sur
$\hat{A}$ dans cette formule. Le reste : « onto » pour « unto »,
« commutes with », « standard », « developed », les contractions
(« isn't », « doesn't »), « a long time ago ».}

\page{22}
\note{p.\ 156 à 190 du tapuscrit. Surtout l'orthographe
(« characterisation », « developing », « squarable » mis en doute) et
une règle générale : \emph{non} prend un trait d'union
(« non-empty », « non-abelian »). À la p.\ 167 l'auteur de la liste réécrit
une formule \uncertain{$f^\circ \colon C^\circ \to C'^\circ$}. À la
p.\ 186, deux phrases jugées fautives, dont une sur la définition des
« intervalles homotopiques » non faibles, qu'il propose de tourner
autrement. À la p.\ 188, une fin de ligne manque sur sa copie ; à la
p.\ 190, « manifold » et non « manyfold ».}

\page{23}
\note{p.\ 191 à 224 du tapuscrit. « recense » n'est pas anglais :
« check-on » est proposé ; « datum » et non « data » au singulier
(p.\ 207) ; des phrases à reprendre (p.\ 193, 210). À la p.\ 216, il
demande si l'on veut dire « that the functor », la ligne n'ayant pas de sens
en l'état. À la p.\ 224, en renvoyant à la p.\ 32 :
$j_A \colon \mathrm{Cat} \to \hat{A}$,
$j_A(X) = \bigl(a \mapsto \mathrm{Hom}(A_{/a}, X)\bigr)$, et une
ligne du tapuscrit est jugée fausse en conséquence (quelques mots
interlinéaires, \ill{}).}

\page{24}
\note{p.\ 225 à 274 du tapuscrit. À la p.\ 234, une reformulation
proposée d'un énoncé sur $W$ et $(A, i)$ ; à la p.\ 235, l'énoncé
« $x$ asphérique ssi $a$ asphérique » est relevé comme une tautologie
telle qu'il est écrit. À la p.\ 240,
$W = (i_{M,i})^{-1}(\underline{W})$ ; à la p.\ 246, une question en
interligne sur la condition (ii) : ne faut-il pas
$i_0 \colon N_0 \to N$ ? À la p.\ 248, un ajout manuscrit du tapuscrit
relu « tel que $\hat{C}$ est totalement… » ; à la p.\ 251,
$i \colon A \to M$. À la p.\ 269, $\pi_0$, $\pi_1$ ne sont pas définis à
cet endroit et la notation est dite confuse. Le reste est d'usage
(contractions, « its », « developed »).}

\page{25}
\note{p.\ 275 (le numéro n'est pas porté) à 302 du tapuscrit. En tête,
une question sur \uncertain{« $C \simeq A$ ? »}, jugée confuse en l'état.
À la p.\ 288, une phrase sans verbe (« hence … $M'$ ») est à compléter.
« commutes with » et « commutation … with » reviennent à presque chaque
page ; « analogons » n'est pas au dictionnaire, « analogues » est
proposé (p.\ 296) ; une tournure à revoir à la p.\ 293.}

\page{26}
\note{p.\ 302 à 334 du tapuscrit. À la p.\ 303, « image of $A$ by $A'$ »
n'est pas compris. À la p.\ 315, une longue phrase jugée difficile, dont
l'auteur de la liste propose une réécriture (sur ce qui, dans l'état
présent des connaissances, distingue une catégorie pour exprimer les
situations de base). À la p.\ 327, faut-il lire
$\alpha \colon \Delta \to \Delta_{/F}$ ? À la p.\ 329, un $\Delta$ dont
le trait doublé manque sur la copie, et « semisimplicial » ; en marge
droite, un doute sur « enveloping ». À la p.\ 334,
« $b \in A_{n-1}$ ».}

\page{27}
\note{p.\ 335 à 354 du tapuscrit. Copie, dans un état antérieur, des
vingt premières lignes de la page 28 : mêmes renvois et mêmes corrections,
sans la date ni la suite. Rien que d'usage (« skilful », « afraid »,
« commutes with », « independent of », « inadequacies »).}

\page{28}
\note{p.\ 335 à 371 du tapuscrit. Après la p.\ 354, un trait et la mention
« Sent 6/2/84 », puis la liste reprend à la p.\ 354. À la p.\ 355,
l'auteur demande ce que l'on veut dire « from $(*)$ ». Pour le reste,
contractions et « commutes with ».}

\page{29}
\note{p.\ 371 à 401 du tapuscrit. À la p.\ 385, le signe $\simeq$ manque
dans la formule (7) ; à la p.\ 401, « $(k\text{-}M)$ ». « with » et non
« to » après « commute » (p.\ 379) ; une faute de frappe (p.\ 382) ;
« interpret », « made up of », « inaccurate ».}

\page{30}
\note{p.\ 402 à 425 du tapuscrit. À la p.\ 402, l'indice de
$\hat{A}_{\mathrm{ab}}$ manque sur la copie. À la p.\ 407, une remarque
sur « formulaire » : le mot anglais est « formulary », jugé excellent,
et il revient aux p.\ 411 à 422 ; « relying » n'est pas employé en ce sens,
« binding » est proposé. Aux p.\ 408 et 411, des lettres manquent sur la
copie ; à la p.\ 409, un passage n'est pas compris ; aux p.\ 413 et 417,
des renvois aux formules (75) et (84).}

\page{31}
\note{p.\ 426 à 468 du tapuscrit. À la p.\ 428, « $\hat{A}$ is abelian »
est relevé comme une erreur : c'est $\hat{A}_M$ qui est voulu. À la
p.\ 431, « integration » n'est pas clair ; à la p.\ 438, une parenthèse
manque dans une formule ; aux p.\ 440 et 449, deux phrases reformulées
(sur la flèche diagonale, et sur des groupes d'homotopie qui ne sont pas des
$k$-modules). « quasi-isomorphisms », « divided », « mimic »,
« Postnikov's ».}

\page{32}
\note{p.\ 468 à 482 du tapuscrit. À la p.\ 471, une phrase jugée ambiguë
est récrite : un complexe $X_*$, lisse et satisfaisant la condition de Kan,
devrait donner un objet en groupes $G = \pi_1(X_*)$ dans $U(k)$, avec une
submersion (H) ; l'auteur de la liste ajoute qu'il a pu se méprendre sur le
sens voulu. À la p.\ 476, la virgule autour d'une incise (« la seule
condition à imposer ») ; à la p.\ 481, une phrase sur la formule (136) à
reprendre. Le reste d'usage.}

\page{33}
\note{p.\ 483 à 494 du tapuscrit. À la p.\ 484, une flèche manque dans la
formule (8) ; à la p.\ 485, l'expression
\uncertain{$W^{\wedge}\,\Gamma^{\wedge}_k(M)$} ; à la p.\ 490, des
reformulations d'un passage de récit. À la p.\ 494, une parenthèse est mal
placée dans $\prod_j \Gamma^j_k(M)$ et la ligne est jugée confuse ;
l'auteur de la liste demande une explication explicite des notations
$\mathrm{Sym}^j_k$ et $\Gamma^j_k$, dont il a dû chercher le sens ailleurs
(SGA 6 pour $\Gamma_k$). Ligne 7, $\Gamma^{\wedge}_k(M)$ ; en bas,
« considered » a perdu sa fin.}

\page{34}
\note{p.\ 495 à 514 du tapuscrit. À la p.\ 495, l'auteur de la liste
demande si « endomorphism » est juste là, les groupes d'endomorphismes étant
rares, en avouant mal connaître les schémas en groupes additifs. À la
p.\ 496, un « 0 » manque dans une cohomologie ; à la p.\ 499, un mot est
trop pâle pour être lu sur la copie ; à la p.\ 506, « Lazare » ou
« Lazard » ? À la p.\ 509, il présume « between $LH$ and $E$ », et
note que la formule (8) se factorise par (2$'$), non par (2).
« dévissage », « Whitehead's ».}

\page{35}
\note{p.\ 515 à 532 du tapuscrit. À la p.\ 516, un doute sur « let's »
et « lets », qu'il lui demande de vérifier. À la p.\ 517, que désigne
$V'$, puisque $V$ et $V'$ sont tous deux employés ? et « of Breen ». À la
p.\ 529, $X^{\vee}$ ; aux p.\ 527--528, un « ons » non supprimé et une
phrase à reprendre. En bas, une ligne biffée.}

\page{36}
\note{p.\ 532 à 549 du tapuscrit. À la p.\ 538, la formule (6) paraît
étrange, l'indice de $\Pi_i$ étant placé ailleurs ; l'auteur de la liste
propose plusieurs lectures sans trancher : $H^i(T, \Pi_i(X_*)) = 0$ pour
$i > 0$, ou $H^i(T, \Pi_j(X_*)) = 0$ pour $j > 0$, ou pour $i \geq j$,
$i, j > 0$. « of Illusie », « J.H.C. Whitehead », « dissymmetry » ;
aux p.\ 543 et 547, deux tournures à revoir.}

\page{37}
\note{p.\ 550 à 557 du tapuscrit, dernière page de la liste, qui
s'arrête au tiers de la feuille. « formulary » à plusieurs lignes de la
p.\ 550, « Hurewicz », et des corrections d'usage dans un passage de récit
(p.\ 555--557).}

\end{document}
